\documentclass[conference]{IEEEtran}
\usepackage{amsmath}
\usepackage{enumitem}

\title{Frozen Paper Specification: Stage-Wise Audit of a TDA Poisoning Detector}

\author{\IEEEauthorblockN{Christian Dane Beels}
\IEEEauthorblockA{Department of Mathematical Sciences\\
United States Military Academy\\
West Point, New York}}

\begin{document}
\maketitle

\section*{Status and Purpose}
\textbf{Frozen August 2, 2026.} This document fixes the scope, comparison,
primary endpoint, and permissible claims for a five-page MIT URTC paper before
any additional experiment is considered. The object of study is a
\emph{controlled reconstruction} of the published TDA--clustering approach,
not ``the Monkam detector'' and not an exact reproduction of the source study.

\section*{Research Question}
\textbf{Where does information needed to distinguish byte-permutation
poisoning disappear in a controlled reconstruction of a TDA-based
network-traffic detector, and does eliminating single-threshold representation
collisions improve downstream poison removal?}

\section*{Frozen Experimental Design}
\begin{itemize}[leftmargin=*,nosep]
  \item \textbf{Dataset:} only the packet-level UNSW-NB15 Payload-Byte
  derivative: 1,500 padded payload-byte fields per observation. Each seeded
  arm uses the same 5,000-row clean sample and 500 appended poisoned rows.
  Seeds are \(\{42,123,456,789,1024\}\).
  \item \textbf{Attack families:} support-restricted, guaranteed-nontrivial
  permutations: 60 disjoint transpositions; one 120-byte block reversal; one
  swap of two disjoint 60-byte blocks; and one nonzero cyclic shift of the
  conservative-support prefix. Claims are conditional on these four
  prespecified families.
  \item \textbf{Control:} the legacy threshold-0.4 representation and its
  60-dimensional feature vector.
  \item \textbf{Only intervention:} the fixed threshold stack
  \(\{0.1,0.2,\ldots,0.9\}\), producing nine 60-feature blocks concatenated in
  threshold order and divided by \(\sqrt{9}\), for 540 features. The same
  poisoned realization is used in both arms.
  \item \textbf{Held fixed:} 30-by-50 raster; two Height and three Radial
  filtrations; cubical persistence; Scaler; persistence entropy and five
  amplitude summaries; and \(\mathrm{OPTICS}(\texttt{min\_samples}=5,
  \texttt{max\_eps}=2.0)\).
\end{itemize}

\section*{Frozen Evidence and Primary Endpoint}
Q3's mutually exclusive earliest-merger attribution is evidence of record.
Across five seeds, \textbf{50.80\% $\pm$ 0.50\%} of final repeated-member
collision mass first merges as identical raw 1,500-byte payloads,
\textbf{46.68\% $\pm$ 0.46\%} first merges at threshold-0.4 binarization, and
\textbf{2.52\% $\pm$ 0.31\%} first merges in cubical persistence. The support
record, five filtration images, Scaler, persistence summaries, and final
concatenation create no additional mergers. These percentages partition
collision mass; they are not percentages of all 5,500 rows.

The \textbf{primary detector endpoint} is the repair-minus-control difference
in poison-removal rate at matched clean false-removal cost. At each control
point with retrospective cluster-purity threshold
\(\{1.0,>0.95,>0.90,>0.80,>0.50\}\), the repair is evaluated at its best
prespecified point whose clean false-removal rate does not exceed the control
budget. Exact-purity capture, removal precision, poison/clean unclustered
fractions, cluster count, exact collisions, and displacement diagnostics are
secondary. Ground-truth cluster purity is an evaluation instrument, not a
deployable cluster-labeling rule.

\section*{Claim Boundary and Analysis Lock}
The paper may distinguish algebraic representation repair, statistical
separation, and downstream detection; success at one level does not imply
success at another. It may report a controlled negative result. It may not
claim robustness to arbitrary permutations, general TDA failure, exact
reproduction of the source study, or detector improvement from collision
reduction alone.

\textbf{No post hoc tuning is permitted.} After observing the recorded results,
the paper will not select thresholds, learn or change threshold-block weights,
add/drop/reweight filtrations, alter vector summaries, tune OPTICS or its
cluster-extraction rule, resample attacks, delete seeds, or add a second dataset
to improve the reported outcome. Any later mechanism test must be separately
preregistered, analyzed as a new experiment, and excluded from the frozen
comparison unless explicitly approved before it is run.

\end{document}
